\section{Related Work}
\label{sec:related}

\paragraph{Security framework crosswalks.}
Manual crosswalks between cybersecurity frameworks have a long history.
NIST maintains mappings between SP~800-53 and ISO~27001, and the Cloud
Security Alliance publishes the Cloud Controls Matrix with pre-built
mappings to multiple standards. In the AI-specific domain, the
\emph{GenAI Security Project} has produced a community-driven crosswalk
covering several major AI security
frameworks~\citep{genai_security_project}. Our work uses the GenAI
Security Project crosswalk as an upstream data source for prior-edge
construction and acknowledges its contribution under the CC~BY-SA~4.0
license. To our knowledge, no prior work has attempted to
\emph{automate} the crosswalk process for AI security frameworks using
machine learning at this scale.

\paragraph{Graph neural networks for knowledge graphs.}
Graph attention networks~\citep{velickovic2018gat} and their successors
such as GATv2~\citep{brody2022gatv2} have been widely used for link
prediction and node classification in knowledge graphs. While our
current feature pipeline extracts structural graph features directly
rather than training a GATv2 end-to-end, the graph construction and
feature extraction approach is informed by this literature, and
GATv2-based node embeddings remain a planned extension.

\paragraph{Conformal prediction for classification.}
Conformal prediction~\citep{vovk2005conformal} provides
distribution-free coverage guarantees for classification. Mondrian
conformal prediction~\citep{romano2020malice} extends this to
per-class calibration, which is essential in our setting due to severe
class imbalance across the four tiers. We adopt Mondrian conformal
prediction to provide per-tier coverage guarantees for the crosswalk
classifier.

\paragraph{LLM-as-judge evaluation.}
Recent work has demonstrated that large language models can serve as
effective evaluators for subjective tasks, including semantic similarity
judgment and text classification. We incorporate LLM-as-judge tier
scores as features in our stacking classifier, treating the LLM
judgment as a calibrated prior that the gradient boosting model refines
using structural and embedding-based evidence.

\paragraph{Gradient-boosted ensembles.}
LightGBM~\citep{ke2017lightgbm} is a widely adopted gradient boosting
framework that handles heterogeneous feature types (dense embeddings,
sparse lexical scores, graph metrics, and discrete LLM judgments) in a
single model. We use LightGBM as our stacking classifier over features
from all three signal families.
